\documentclass[10pt,a4paper]{amsart}
\usepackage[margin=19mm]{geometry}
\usepackage{amsmath,amssymb,mathtools}
\usepackage[scr=rsfs,cal=euler]{mathalfa}
\usepackage{xfrac}
\usepackage[unicode,hidelinks,bookmarks=false]{hyperref}
\newcommand{\ie}{i.e.\ }
\newcommand{\cf}{cf.\ }
\newcommand{\resp}{respectively\ }
\newcommand{\Subsection}{\subsection}
\providecommand{\nobreakdash}{\nobreak\hbox{-}}
\setlength{\emergencystretch}{2em}
\multlinegap0pt
\usepackage{amssymb}
\usepackage{enumerate}
\usepackage{tikz}
\usepackage{mathtools}
\newtheorem{proposition}{Proposition}
\title{A comparison of the v-number of a monomial ideal and its integral closure}
\author{Prativa Biswas, Mousumi Mandal, Partha Phukan}
\date{Source: arXiv:2609.05044v1 (CC BY 4.0)}
\begin{document}
\maketitle

\noindent\emph{Source and licence.} A comparison of the v-number of a monomial ideal and its integral closure. Version arXiv:2609.05044v1 (CC BY 4.0), distributed by arXiv under the Creative Commons Attribution 4.0 International licence, http://creativecommons.org/licenses/by/4.0/, as recorded in the arXiv licence metadata for this exact version and shown on the abstract page https://arxiv.org/abs/2609.05044v1. Full licence notice: \texttt{licenses/arxiv-2609.05044-CC-BY-4.0.txt}.\par\smallskip

\noindent\textit{Editorial scope.} Counted unit: the article's proposition prop:v-stability-squarefree (source0.tex lines 1211--1218). The article's complete original proof of it is delivered with it (source0.tex lines 1220--1275, one \texttt{proof} environment of 56 lines). No mathematics is written, repaired or re-derived: the statement and the proof are the article's own lines, in the article's own notation, and no retained line is substituted or re-flowed.  No further source-local context is needed: every cross-reference of every retained line resolves inside the two delivered spans.  Nothing was omitted from any retained span, so the package carries no omission record.  No retained line contains a citation, so the wrapper carries no bibliography and no bibliography entry.  No retained line contains a figure environment, an image-inclusion command or drawing code, so no image is delivered and the package declares no asset dependency.  Editorial formatting only, in addition: an AMS \texttt{amsart} preamble that restates the article's own declarations and macros used by the retained lines, namely the environments \texttt{proposition} on separate counters, together with the standard packages those lines need.  The retained environments print in this document as Proposition 1 in order of appearance, whereas the article prints the same items with the numbering of its own full document.  The complete native source of the article as distributed in the arXiv e-print for this exact version is bundled unchanged as \texttt{sources/209410/source0.tex}.
\begin{proposition}
\label{prop:v-stability-squarefree}
Let \(I\) be a squarefree monomial ideal such that \(\mathrm{v}(I) = \alpha(I) - 1\). Then
\[
\mathrm{v}(\overline{I^k}) = \mathrm{v}(I^k) = k\alpha(I) - 1
\]
for all \(k \ge 1\).
\end{proposition}

\begin{proof}
Let \(f\) be a squarefree monomial such that
\(
\deg f = \alpha(I) - 1 \quad \text{and} \quad (I : f) = \mathfrak{p}.
\)

\noindent Then \(x_{i_j}\) does not divide \(f\) for all \(x_{i_j} \in \mathcal{G}(\mathfrak{p})\). Let \(\mathfrak{p} = (x_{i_1}, \ldots, x_{i_s})\).
\noindent We claim that
\[
(\overline{I^k} : x_{i_1}^{k-1} f^k) = \mathfrak{p}.
\]
Indeed, since
\(
x_{i_j} x_{i_1}^{k-1} f^k = (x_{i_1}^{k-1} f^{k-1})(x_{i_j} f) \in I^k \subseteq \overline{I^k},
\)
it follows that \(\mathfrak{p} \subseteq (\overline{I^k} : x_{i_1}^{k-1} f^k)\).

\noindent In order to prove the other inclusion, let \(u = x_1^{a_1} \cdots x_n^{a_n}\) be a minimal generator of \(I^k\). Define
\(
w_{\mathfrak{p}}(u) = {\sum_{x_j \in \mathfrak{p}}} a_j.
\)
Then \(w_{\mathfrak{p}}(u) \ge k\) for all \(u \in \mathcal{G}(I^k)\).
All exponent vectors of monomials in \(I^k\) lie in 
\(
\left\{ \mathbf{a} : \sum_{x_j \in \mathfrak{p}} a_j \ge k \right\}.
\)
Hence, every point in \(NP(I^k)\) lies in
\(
\left\{ \mathbf{a} : \sum_{x_j \in \mathfrak{p}} a_j \ge k \right\}.
\)
Note that, \(w_{\mathfrak{p}}(f) = 0\), so
\(
w_{\mathfrak{p}}(x_{i_1}^{k-1} f^k) = k - 1.
\)
For any monomial \(v \notin \mathfrak{p}\), we have \(w_{\mathfrak{p}}(v) = 0\), and therefore
\(
w_{\mathfrak{p}}(v x_{i_1}^{k-1} f^k) = k - 1.
\)
Thus, \(v x_{i_1}^{k-1} f^k \notin \overline{I^k}\), which implies
\(
v \notin (\overline{I^k} : x_{i_1}^{k-1} f^k).
\)
Hence,
\[
(\overline{I^k} : x_{i_1}^{k-1} f^k) = \mathfrak{p}.
\]
Therefore,
\[
\mathrm{v}(\overline{I^k}) \le \deg(x_{i_1}^{k-1} f^k) = (k-1) + \deg f^k = (k-1) + k(\alpha(I) - 1) = k\alpha(I) - 1.
\]
Since \(\mathrm{v}(I^k) \ge k\alpha(I) - 1\) for all \(k \ge 1\), we obtain
\[
\mathrm{v}(\overline{I^k}) = \mathrm{v}(I^k) = k\alpha(I) - 1.
\]
This completes the proof.
\end{proof}

\end{document}
